% Options for packages loaded elsewhere
\PassOptionsToPackage{unicode}{hyperref}
\PassOptionsToPackage{hyphens}{url}
\documentclass[
  ignorenonframetext,
]{beamer}
\newif\ifbibliography
\usepackage{pgfpages}
\setbeamertemplate{caption}[numbered]
\setbeamertemplate{caption label separator}{: }
\setbeamercolor{caption name}{fg=normal text.fg}
\beamertemplatenavigationsymbolsempty
% remove section numbering
\setbeamertemplate{part page}{
  \centering
  \begin{beamercolorbox}[sep=16pt,center]{part title}
    \usebeamerfont{part title}\insertpart\par
  \end{beamercolorbox}
}
\setbeamertemplate{section page}{
  \centering
  \begin{beamercolorbox}[sep=12pt,center]{section title}
    \usebeamerfont{section title}\insertsection\par
  \end{beamercolorbox}
}
\setbeamertemplate{subsection page}{
  \centering
  \begin{beamercolorbox}[sep=8pt,center]{subsection title}
    \usebeamerfont{subsection title}\insertsubsection\par
  \end{beamercolorbox}
}
% Prevent slide breaks in the middle of a paragraph
\widowpenalties 1 10000
\raggedbottom
\AtBeginPart{
  \frame{\partpage}
}
\AtBeginSection{
  \ifbibliography
  \else
    \frame{\sectionpage}
  \fi
}
\AtBeginSubsection{
  \frame{\subsectionpage}
}
\usepackage{iftex}
\ifPDFTeX
  \usepackage[T1]{fontenc}
  \usepackage[utf8]{inputenc}
  \usepackage{textcomp} % provide euro and other symbols
\else % if luatex or xetex
  \usepackage{unicode-math} % this also loads fontspec
  \defaultfontfeatures{Scale=MatchLowercase}
  \defaultfontfeatures[\rmfamily]{Ligatures=TeX,Scale=1}
\fi
\usepackage{lmodern}
\ifPDFTeX\else
  % xetex/luatex font selection
\fi
% Use upquote if available, for straight quotes in verbatim environments
\IfFileExists{upquote.sty}{\usepackage{upquote}}{}
\IfFileExists{microtype.sty}{% use microtype if available
  \usepackage[]{microtype}
  \UseMicrotypeSet[protrusion]{basicmath} % disable protrusion for tt fonts
}{}
\makeatletter
\@ifundefined{KOMAClassName}{% if non-KOMA class
  \IfFileExists{parskip.sty}{%
    \usepackage{parskip}
  }{% else
    \setlength{\parindent}{0pt}
    \setlength{\parskip}{6pt plus 2pt minus 1pt}}
}{% if KOMA class
  \KOMAoptions{parskip=half}}
\makeatother
\setlength{\emergencystretch}{3em} % prevent overfull lines
\providecommand{\tightlist}{%
  \setlength{\itemsep}{0pt}\setlength{\parskip}{0pt}}
% Slide layout defaults for warning-clean scientific decks.
\usepackage{etoolbox}
\IfFileExists{xurl.sty}{\usepackage{xurl}}{}
\IfFileExists{seqsplit.sty}{\usepackage{seqsplit}}{\newcommand{\seqsplit}[1]{#1}}
\protected\def\breakseq#1{\seqsplit{#1}}
\protected\def\breaktt#1{\begingroup\ttfamily\seqsplit{#1}\endgroup}
\setlength{\emergencystretch}{6em}
\tolerance=5000
\hbadness=10000
\hfuzz=1pt
\setlength{\tabcolsep}{2pt}
\AtBeginEnvironment{longtable}{\tiny\renewcommand{\arraystretch}{0.86}\setlength{\tabcolsep}{1pt}}
\AtBeginEnvironment{tabular}{\tiny\renewcommand{\arraystretch}{0.86}\setlength{\tabcolsep}{1pt}}

% Natbib and cross-reference fallbacks — slides don't load natbib
% or cleveref, but combined-PDF manuscript prose may emit these
% commands. The fallback renders citations as a bracketed key list
% and cross-references as detokenized labels so slides don't fail on
% undefined control sequences or raw underscores. \providecommand is
% a no-op if packages load later.
\providecommand{\citep}[1]{[#1]}
\providecommand{\citet}[1]{#1}
\providecommand{\citealp}[1]{#1}
\providecommand{\citeauthor}[1]{#1}
\providecommand{\citeyear}[1]{#1}
\providecommand{\cref}[1]{\texttt{\detokenize{#1}}}
\providecommand{\Cref}[1]{\texttt{\detokenize{#1}}}
\usepackage{bookmark}
\IfFileExists{xurl.sty}{\usepackage{xurl}}{} % add URL line breaks if available
\urlstyle{same}
% fallback for those not using the hyperref driver hyperxmp:
\makeatletter
\@ifundefined{xmpquote}{\newcommand{\xmpquote}[1]{#1}}{}
\makeatother
\hypersetup{
  hidelinks,
  pdfcreator={LaTeX via pandoc}}

\author{\texorpdfstring{}{}}
\date{}

\begin{document}

\begin{frame}[fragile,allowframebreaks]{Limitations}
\protect\phantomsection\label{limitations}
This section states what the exemplar does \emph{not} do, alongside what
it does. Coverage and test-count figures are not restated as literal
numbers here --- those live only in the \texttt{493} / \texttt{96.28}
tokens, resolved at render time from a live measurement
(\breaktt{src/manuscript\_variables.py::measure\_test\_summary}), not
hand-typed.

\begin{block}{Reserved slots are excluded from the effective product
space}
\protect\phantomsection\label{reserved-slots-are-excluded-from-the-effective-product-space}
The grammar declares six slots, but \texttt{materialize()} only branches
on three of them --- \breaktt{primitive\_domain}, \texttt{track},
\texttt{section\_set}. The remaining 3
(\breaktt{figure\_profile,\ qr\_profile,\ integrity\_profile}) are parsed
and contribute to \texttt{spec\_hash} like any other slot, but
\breaktt{manuscript\_figures.py}, \texttt{sealing.py}, and
\texttt{integrity.py} are unconditionally exercised regardless of
\texttt{figure\_profile}, \texttt{qr\_profile}, or
\breaktt{integrity\_profile} --- those modules do not currently branch on
the reserved slots' selected option. The practical consequence: distinct
seeds can produce a nominally distinct \texttt{spec\_hash} while
materializing byte-identical children, since only three slots vary
output. That inflation, 360 nominal vs.~45 effective, is disclosed
rather than hidden: \breaktt{generate\_variables}
(\breaktt{src/manuscript\_variables.py}) exposes both
\texttt{PRODUCT\_SIZE} and \breaktt{EFFECTIVE\_PRODUCT\_SIZE} as separate
tokens, so nowhere in this manuscript can the larger, nominal number be
quoted without the smaller, effective one appearing beside it. Wiring
the reserved slots into \texttt{materialize()} is an open item, not a
claimed capability.
\end{block}

\begin{block}{No child-PDF rendering in CI}
\protect\phantomsection\label{no-child-pdf-rendering-in-ci}
\breaktt{realize.py::render\_child\_manuscript()} writes manuscript stubs
for a generated child but does not invoke Pandoc or Chrome. Absent a
full rendering toolchain --- the default posture in CI --- it returns
\texttt{\{"success":\ False,\ "reason":\ "Chrome/Pandoc\ rendering\ not\ available\ in\ child\ context",\ ...\}}
rather than raising or silently skipping. Verifying a child's
\emph{files} (\texttt{verify\_child} / \breaktt{verify\_child\_full}) is
decoupled from verifying that its manuscript renders; only the former is
tested here.
\end{block}

\begin{block}{Within-platform guarantee only}
\protect\phantomsection\label{within-platform-guarantee-only}
Determinism is derived from \texttt{hashlib.sha256} over Python strings
(\texttt{expand.py}'s seed/slot/ordinal/options digest, and
\breaktt{integrity.py::tree\_hash\_from\_content\_hashes} over file
contents), which is itself platform-independent. But the surrounding
pipeline --- file iteration order ahead of \texttt{sorted()}
normalization, path-separator handling, the specific CPython minor
version running \texttt{materialize()} --- is not independently fuzzed
across platforms in this suite. What is actually tested is determinism
within one CPython version on one OS, a narrower claim than the
cross-platform, bit-for-bit reproducibility discussed in the
reproducible-builds literature \breakseq{[@reproducible\_builds]}. A second CI
leg that materializes the same seed on a second OS/Python and diffs tree
hashes would close this gap; it does not currently exist.
\end{block}

\begin{block}{Grammar does not self-modify}
\protect\phantomsection\label{grammar-does-not-self-modify}
The autopoiesis metaphor is figurative. Maturana and Varela's
operational sense of autopoiesis is a system that continuously
regenerates its own constitutive components through its own operation
\breakseq{[@maturana\_varela\_1980]}. Here the grammar
(\breaktt{manuscript/config.yaml}) is fixed input;
\texttt{parse\_grammar} and \texttt{expand} are pure functions of that
input plus a seed; no code path feeds a materialized child back into the
grammar or rewrites \texttt{src/grammar.py}. Children are causally
downstream of the grammar --- the reverse direction does not occur in
the current codebase. The name is a provocation about what genuine
self-production would require, not a claim this exemplar achieves it.
\end{block}

\begin{block}{Integrity verification checks self-consistency, not
tamper-proofing}
\protect\phantomsection\label{integrity-verification-checks-self-consistency-not-tamper-proofing}
\texttt{verify\_child} recomputes a tree hash from the files
\emph{listed inside} \texttt{provenance.json} and compares it against
the \texttt{tree\_hash} field stored in that same file
(\texttt{src/verify.py}, \breaktt{src/materialize.py}). Both the manifest
and the expected hash are self-reported at materialization time; nothing
external anchors them. An actor who can rewrite \texttt{provenance.json}
can edit its \texttt{files} list and recompute a matching hash from
whatever content they choose, and \texttt{verify\_child} would report a
match --- it only checks internal consistency, not consistency against
something outside the child's own directory
\breakseq{[@merkle\_tree\_provenance]}. \texttt{verify\_seal} similarly checks
the shape of an optional \texttt{seal.json} written by the same run it
would need to audit. This detects accidental post-generation drift, the
tested use case; it is not a defense against deliberate tampering with
write access. Also worth naming precisely:
\breaktt{tree\_hash\_from\_content\_hashes} sorts
\breaktt{path:content\_hash} pairs and hashes their concatenation once
--- a flat manifest digest, not a hierarchical Merkle tree with per-file
inclusion proofs. \texttt{integrity.py} separately defines
\texttt{merkle\_root}, an actual pairwise-reduction tree, but
\texttt{materialize()}/\texttt{verify()} call the flat function, not
that one.
\end{block}

\begin{block}{Coverage is uneven across modules, not uniform (though
every module now clears the floor)}
\protect\phantomsection\label{coverage-is-uneven-across-modules-not-uniform-though-every-module-now-clears-the-floor}
Two earlier drafts of this section each named a different below-floor
trio --- first \texttt{sealing.py}/\texttt{verify.py}/\texttt{cli.py},
then, after those were hardened,
\texttt{common.py}/\texttt{figures.py}/\texttt{cover\_art.py}. As of
this measurement no module sits below the 90\% branch-coverage line
(\breakseq{[@fig:coverage\_by\_module]}): dedicated tests were added for
\texttt{common.py}'s \texttt{trunc()} clipping branch and
\breaktt{CheckReport.failed} (\breaktt{tests/test\_common.py}, new this
session), \texttt{figures.py}'s \texttt{list}/\texttt{tuple} input
branch of \breaktt{\_first\_plottable\_array}, the generic
\texttt{repr()} fallback in \breaktt{\_scalar\_summary\_lines}, and the
array-plotting branch of \breaktt{render\_primitive\_figure}
(\breaktt{tests/test\_figures.py}), and \texttt{cover\_art.py}'s QR-seal
drawing branch (\breaktt{tests/test\_cover\_art.py},
\breaktt{test\_render\_cover\_with\_grammar\_hash\_*}) --- the same
branch identified above as running in production but previously
untested.

Smaller residual gaps remain in modules that are otherwise well above
the floor: \texttt{sealing.py}'s \texttt{read\_qr\_matrix}
\texttt{pyzbar}-decode path is untested since \texttt{pyzbar} is not
installed here; \texttt{verify.py}'s schema-version and malformed-JSON
\texttt{except} branches are not independently exercised;
\texttt{cli.py}'s \texttt{cmd\_materialize}/\texttt{cmd\_verify} have no
end-to-end CLI-level test, though the library functions they wrap are
covered elsewhere in \breaktt{test\_materialize.py} and
\breaktt{test\_integrity\_and\_verify.py}; \texttt{honesty.py} itself
sits closest to the floor at just above 90\%. None of these pull the
aggregate below the 90\% floor (\breaktt{-\/-cov-fail-under=90}), and the
aggregate 96.28 alone would obscure exactly where the remaining, smaller
gaps live --- which is the reason this section, and
\breakseq{[@fig:coverage\_by\_module]}, exist as a per-module view rather than
trusting one headline number. Property-based tests
(\breaktt{tests/test\_property\_invariants.py}, Hypothesis) and a
stress/edge suite (\breaktt{tests/test\_stress\_edge\_cases.py}) cover
invariants like boundary seeds and all-reserved-slot configurations
\breakseq{[@claessen2000quickcheck; @maciver2019hypothesis]}, but generated
inputs are not a substitute for direct tests of the specific branches
named above.

This whole section is itself evidence for a point made elsewhere in this
manuscript: a coverage ranking is a measurement, not a fact about the
code --- each of its two prior versions was accurate when written and
became false as soon as new tests shipped. Readers should treat any
specific module name and percentage in this document as dated to its
stated measurement, and re-run \breaktt{stage\_02\_analysis.py} for the
current ranking rather than trust a prior render.
\end{block}

\begin{block}{Cover art now includes the QR seal, but that code path is
untested}
\protect\phantomsection\label{cover-art-now-includes-the-qr-seal-but-that-code-path-is-untested}
An earlier draft of this section reported that the title-page cover
image omitted the originally envisioned QR seal, gradient glow, and
seed-derived dot placement. That is no longer accurate and is corrected
here rather than left to silently drift:
\breaktt{scripts/generate\_cover\_art.py} now calls
\breaktt{render\_cover(...,\ grammar\_hash=grammar.grammar\_hash)}, and
\breaktt{src/cover\_art.py} draws all three elements unconditionally
except the QR seal, which is drawn whenever \texttt{grammar\_hash} is
not \texttt{None} --- the shipped \breaktt{paper.cover.image}
(\breaktt{output/figures/cover\_art.png}) is generated by exactly this
call, so the rendered title page carries a real gradient glow, real
seed-derived dot scatter, and a real QR-style pixel grid encoding
\breaktt{f84a8f9dbcb18e37} with the hash printed as a text label beneath
it. An earlier draft of this section also noted that no test called
\texttt{render\_cover} with an explicit \texttt{grammar\_hash=}, leaving
the QR-drawing branch (\breaktt{src/cover\_art.py}, the
\breaktt{if\ grammar\_hash\ is\ not\ None:} block) exercised in
production but untested; that gap is now closed by
\breaktt{tests/test\_cover\_art.py::test\_render\_cover\_with\_grammar\_hash\_*},
added this session (see ``Coverage is uneven across modules'' below).

Separately, for most of this project's life
\breaktt{manuscript/references.bib} held a single self-referential
\breaktt{friedman2026autopoiesis} entry noting the DOI was forthcoming,
while \breaktt{99\_references.md} carried a hand-written, uncited list
alongside it --- a citation without a resolvable BibTeX entry is exactly
the kind of unverifiable claim this project's honesty contract exists to
catch. That gap is now closed twice over: \texttt{references.bib}
carries five real, live-verified external citations, and following this
project's own Zenodo deposit the self-citation's DOI
(\breaktt{10.5281/zenodo.21227869}, verified by a live \texttt{curl}
against \texttt{doi.org} before being written into the entry) is real
rather than forthcoming. \breaktt{99\_references.md} annotates each entry
against the section that relies on it rather than listing citations
independent of the bibliography.
\end{block}
\end{frame}

\end{document}
